\section{采样、混叠与可恢复性}
\label{sec:08-Numerical-Analysis/08-Fourier-and-Spectral-Methods/02}

一台设备的振动数据以 \(100\) Hz 采集，谱图上有一个干净的 \(40\) Hz 峰。按图索骥去找转速、找齿数、找轴承的特征频率，凑不出这个 \(40\)。后来才发现，机上有一路 \(60\) Hz 的工频干扰。\(100-60=40\)：那个峰是 \(60\) Hz 分量在采样之后留下的伪装。

麻烦的地方在于，这件事没法从数据里查出来。手里那串数字，与一个真实存在 \(40\) Hz 分量的信号所产生的数字，逐位相同。上一节讲的连续 Fourier 变换默认能看到整条时间轴上的函数，而实际拿到的是一列样本；这一步交出去的东西，不是精度，是身份。

\subsection{采样把频率轴折叠，折过来的频率不带标签}

以周期 \(T\) 采样一个复指数，把频率挪动采样频率 \(f_s=1/T\) 的整数倍，样本值不变：

\[
e^{2\pi i(\nu+k/T)nT}=e^{2\pi i\nu nT}e^{2\pi ikn}=e^{2\pi i\nu nT},
\qquad k\in\Z.
\]

多出来的因子 \(e^{2\pi ikn}\) 恒等于 \(1\)，因为 \(kn\) 是整数。这一行代数就是全部机制：连续频率轴上每隔 \(f_s\) 取一个点，这一整族频率在采样点上给出同一列数。

\begin{center}
\begin{tikzpicture}[font=\small,>=stealth]
  \begin{scope}[xscale=9.4,yscale=1.25]
    \draw[black!30] (-0.02,0) -- (1.06,0);
    \draw[orange!80!black,thick] plot[domain=0:1.04,samples=261,smooth]
      (\x,{sin(deg(12*pi*\x))});
    \draw[blue!65!black,very thick] plot[domain=0:1.04,samples=161,smooth]
      (\x,{sin(deg(2*pi*\x))});
    \foreach \tn/\vv in {0/0,0.2/0.951,0.4/0.588,0.6/-0.588,0.8/-0.951,1.0/0}{
      \draw[black!35,dashed] (\tn,0) -- (\tn,\vv);
      \node[circle,fill=black,inner sep=1.3pt] at (\tn,\vv) {};
    }
  \end{scope}
  \node[right,font=\footnotesize,blue!55!black] at (9.9,0.55) {\(f_1=1\) Hz};
  \node[right,font=\footnotesize,orange!70!black] at (9.9,-0.55) {\(f_2=6\) Hz};
  \node[font=\footnotesize,black!70] at (4.7,-1.95) {采样率 \(f_s=5\) Hz，样本点共六个};
\end{tikzpicture}

\emph{两条频率相差 \(5\) Hz 的正弦，在每一个采样点上取值完全相同。黑点是全部可得的数据；把连续曲线抹掉之后，没有任何依据把它们分开。}
\end{center}

图里那两条曲线在采样点之外差得很远，可数据只剩下那六个黑点。要从黑点反推曲线，得先假定原信号里没有 \(6\) Hz——这个假定来自采样之外，来自对设备的了解，不来自数据本身。开头那个 \(40\) Hz 峰就是这么来的：\(60\) Hz 与 \(-40\) Hz 相差一个 \(f_s\)，而对实值信号，\(-40\) 与 \(40\) 又共轭配对，于是它现身在 \(40\) Hz 上。

把频率轴按 \(f_s\) 折叠，落在同一位置的所有频率合成一个观测值，这是一次多对一的映射。多对一的映射没有逆。事后做任何数字滤波，能把 \(40\) Hz 处的能量删掉，却答不出这份能量原本来自 \(40\) 还是 \(60\)——那个区分在模数转换的一瞬间就被销毁了，后面的处理链里根本没有它的载体。这是混叠与其他误差的分界：量化噪声、舍入、泄漏都是精度问题，加大字长或延长观测能改善；混叠是信息问题，改善的余地为零。

\subsection{可恢复的条件是两倍带宽，不是一倍}

要让折叠不发生，得让每个折叠格里最多住一个频率。若连续信号严格带限于 \(\abs{\nu}<B\)，且

\[
f_s>2B,
\]

则各个频谱副本互不重叠，中心那一份可以被理想低通原样切出来，再逆变换回时域，得到

\[
x(t)=\sum_{n\in\Z}x(nT)\,\operatorname{sinc}\!\left(\frac{t}{T}-n\right),
\qquad
\operatorname{sinc}(u)=\frac{\sin \pi u}{\pi u}.
\]

推导的骨架只有三步：采样等价于在频域把 \(\hat x\) 以 \(f_s\) 为周期做副本叠加；带宽条件保证副本不重叠；乘上带宽为 \(f_s/2\) 的矩形窗取回中心副本，而理想矩形滤波器的时域响应正是无限支撑 sinc，频域相乘换成时域卷积就得到上式。

条件对账要认真做，因为工程上失守的通常不是不等式而是前提。严格带限：真实信号的谱只是衰减，从不真正为零，超出 \(B\) 的那点尾巴照样折回来。双向无限的样本：重建公式对 \(n\) 求和到 \(\pm\infty\)，截断带来误差，而 sinc 只按 \(1/\abs{t}\) 代数衰减，尾巴掉得很慢，截断误差比想象中大。精确时钟：采样时刻的抖动等价于给信号加了一层随机相位调制，高频分量受害更重。取严格不等号：\(f_s=2B\) 时，恰好落在 Nyquist 频率上的正弦可能被采成全零，相位选得不巧就整条丢掉。

两个名字容易串。给定带宽 \(B\)，\(2B\) 叫 Nyquist rate，是采样率的下界；给定采样率 \(f_s\)，\(f_s/2\) 叫 Nyquist frequency，是可分辨频率的上界。一个是对采样率的要求，一个是对信号的限制，方向相反。

于是实际系统不会把采样率贴着 \(2B\) 设。采样之前先过一道模拟低通，把带外能量压到量化噪声以下；模拟滤波器的过渡带不可能陡直，所以采样率还要为过渡带留出余量，通常取到 \(2.5B\) 以上；剩下的带外 alias 能量要算成一项指标写进设计文档，而不是当作不存在。这些余量都花在采样之前，因为抗混叠只能发生在采样之前。

\subsection{降采样在别处也叫降采样}

同一条道理换个场合就不好认了。

给一张高分辨率照片做缩小，若直接每隔 \(k\) 个像素取一个，图里细密的条纹、格子布、屏幕栅格会变成一片粗大的彩色波纹，也就是摩尔纹。空间频率超过新采样率一半的那些细节没有消失，它们折回到低频，伪装成一种原图里不存在的粗结构。正确的做法是先用低通核（高斯模糊或面积平均）把高频压掉，再抽样——所有像样的图像库都这么做，参数里那个 anti-aliasing 开关就是它。

时间序列同理。把逐日数据抽成每七天一个点，日内和周内的高频成分会折成一个周期几周甚至几个月的``伪周期''，在自相关图上看着像真的。先做滑动平均再抽，和直接抽，得到的两条曲线常常讲两个故事。

卷积网络里的 stride 与池化是同一个操作。带 stride 的卷积等于卷积后抽样，最大池化更是先取局部极值再抽样，两者都没有配套的低通。后果是网络对输入的微小平移变得敏感：图像挪一个像素，抽样格点相对特征的相位就变了，折回来的成分跟着变，输出跟着跳。在抽样前插一层低通（模糊）再抽，平移稳定性明显改善，这正是 anti-aliased downsampling 一类做法的依据。

\subsection{混叠与泄漏要分开记账}

看到频谱不干净，第一件事是判断机制，不能笼统归到``FFT 误差''。混叠来自采样率不足或带外能量，不同的连续频率生成同一个离散频率，不可逆。泄漏来自观察窗有限，单根谱线的能量被窗函数的谱抹到邻近格点上，可以通过换窗在旁瓣高度与主瓣宽度之间权衡，但只要观察时长有限就消不掉。还有一类是真实频率落在两个频率格点之间时峰值幅度被低估，以及幅值离散化引入的量化噪声。四种机制的成因、可逆性、对策都不同，把它们混成一句抱怨，等于放弃了诊断。

其中泄漏与离散频率格点的来历要等下一节讲完 DFT 才好说清，因为它来自 DFT 对数据的一个隐含假设。

最后是这套理论自己的边界。严格时限的信号不可能严格带限，反过来也一样——上一节说过，函数在时域越窄，在频域就越宽，两头都要紧凑是数学上不允许的。所以采样定理描述的是一个理想化情形，实际永远在用足够小的带外能量换足够好的重建。若愿意换掉``带限''这个先验，换成``信号在某组基下稀疏''，非均匀采样与压缩感知能在远低于 Nyquist rate 的采样数下重建。它们不是绕过了采样定理，是往里补充了新的结构假设——被采样销毁的信息，仍然只能靠数据之外的知识补回来。

下一节回到手上那列有限长的样本：把它送进离散 Fourier 变换，看看这个变换到底算的是什么，以及为什么算它可以快得多。
